% !TeX program = xelatex
% Σχολιασμένη βιβλιογραφία -- απαιτεί XeLaTeX (fontspec/polyglossia για ελληνικά).
% Μεταγλώττιση: xelatex bibliography_notes.tex  x2  (από μέσα από docs/)
\documentclass[11pt,a4paper]{article}

\usepackage{fontspec}
\usepackage{polyglossia}
\setmainlanguage{greek}
\setotherlanguage{english}
\setmainfont{Linux Libertine O}
\setsansfont{Linux Biolinum O}
\setmonofont{DejaVu Sans Mono}[Scale=0.85]

\usepackage[a4paper,margin=2.4cm]{geometry}
\usepackage{parskip}
\usepackage{amsmath,amssymb}
\usepackage{enumitem}
\usepackage{xcolor}
\definecolor{linkblue}{HTML}{1a4f8f}
\usepackage[colorlinks=true,urlcolor=linkblue,linkcolor=linkblue]{hyperref}
\setlength{\emergencystretch}{3em}
\newcommand{\en}[1]{\textenglish{#1}}

% Μία πηγή: \src{κλειδί}{πλήρης αναφορά}{σύνδεσμος}
\newcommand{\src}[3]{%
  \par\medskip\noindent\textbf{#2}\\
  \href{#3}{\texttt{\small #3}}\par\smallskip}

\title{\textbf{Σχολιασμένη βιβλιογραφία}\\[6pt]
\large\normalfont Τι χρειάζεται να κρατήσει κανείς από κάθε πηγή,\\
και πού χρησιμοποιείται στη διπλωματική}
\author{}
\date{\today}

\begin{document}
\maketitle

\noindent\small
Κάθε καταχώριση έχει τρία μέρη: την αναφορά, έναν σύνδεσμο, και μια σύνοψη του τι χρειάζεται
πραγματικά από αυτήν --- όχι περίληψη του άρθρου, αλλά το σημείο που ακουμπά την εργασία.
Όπου μια πηγή στηρίζει συγκεκριμένο σημείο του κειμένου, αυτό αναφέρεται ρητά.
\normalsize

% =====================================================================
\section{Θεμέλια της online κυρτής βελτιστοποίησης}
% =====================================================================

\src{zinkevich2003}{Zinkevich, M. (2003). Online convex programming and generalized
infinitesimal gradient ascent. \emph{ICML}, 928--936.}
{https://cdn.aaai.org/ICML/2003/ICML03-120.pdf}
Η εργασία που εισάγει τον \emph{projected OGD} και αποδεικνύει το φράγμα $O(\sqrt{T})$ για
κυρτές απώλειες με φραγμένη κλίση. Κρατάμε τρία πράγματα: τη δομή της απόδειξης με τον
τηλεσκοπικό όρο, το ότι η προβολή σε κυρτό σύνολο \emph{δεν αυξάνει αποστάσεις} (που είναι
ακριβώς το βήμα που κάνει την απόδειξη να δουλεύει), και ότι το βέλτιστο $\eta$ εξαρτάται από
τον ορίζοντα $T$. Το τελευταίο είναι η αφετηρία για όλο το κεφάλαιο συντονισμού βήματος.

\src{hazan2016}{Hazan, E. (2016). Introduction to online convex optimization.
\emph{Foundations and Trends in Optimization}, 2(3--4), 157--325.}
{https://arxiv.org/abs/1909.05207}
Η καλύτερη ενιαία αναφορά για ολόκληρο το πλαίσιο. Αν διαβαστεί ένα μόνο πράγμα για γενική
εποπτεία, αυτό. Περιέχει σε ενιαίο συμβολισμό τα φράγματα OGD, FTRL και exp-concave, καθώς
και το Online Newton Step. Χρήσιμο ιδίως για το κεφάλαιο θεωρίας, όπου οι αποδείξεις
ακολουθούν αυτή τη διατύπωση.

\src{cesabianchi2006}{Cesa-Bianchi, N., \& Lugosi, G. (2006). \emph{Prediction, Learning, and
Games}. Cambridge University Press.}
{https://doi.org/10.1017/CBO9780511546921}
Το βιβλίο αναφοράς για την πρόβλεψη με συμβουλή ειδικών. Από εδώ προέρχονται το λήμμα του
Hoeffding στη μορφή που χρησιμοποιείται στην απόδειξη του Hedge, και --- σημαντικότερο --- το
\emph{κάτω} φράγμα $\Omega(\sqrt{T\ln N})$. Το δεύτερο είναι αυτό που επιτρέπει να ειπωθεί
ότι ο Hedge είναι βέλτιστος ως προς την τάξη μεγέθους και ότι η απόκλιση του συντονισμένου
βήματος δεν σημαίνει αποτυχία της θεωρίας.

\src{littlestone1994}{Littlestone, N., \& Warmuth, M. K. (1994). The weighted majority
algorithm. \emph{Information and Computation}, 108(2), 212--261.}
{https://doi.org/10.1006/inco.1994.1009}
Η πηγή της ιδέας των πολλαπλασιαστικών βαρών. Ιστορικά πρώτη· πρακτικά χρήσιμη για να
φανεί ότι ο Hedge δεν είναι τέχνασμα αλλά η φυσική γενίκευση της σταθμισμένης πλειοψηφίας
σε συνεχή βάρη.

\src{vovk1990}{Vovk, V. G. (1990). Aggregating strategies. \emph{COLT}, 371--386.}
{https://dl.acm.org/doi/10.5555/92571.92672}
Ο Aggregating Algorithm και η έννοια της \emph{αναμιξιμότητας} (\en{mixability}). Το
ουσιώδες για εμάς: για απώλειες που είναι $\alpha$-exp-concave --- και η λογαριθμική είναι,
με $\alpha=1$ --- επιτυγχάνεται regret $\ln N / \alpha$, δηλαδή \emph{ανεξάρτητο του
ορίζοντα}. Η εργασία δείχνει ότι αυτό το θεωρητικά πολύ ισχυρότερο φράγμα συνοδεύεται από
χειρότερη πρακτική επίδοση εδώ, και η εξήγηση είναι ο στατικός comparator.

\src{hazan2007}{Hazan, E., Agarwal, A., \& Kale, S. (2007). Logarithmic regret algorithms for
online convex optimization. \emph{Machine Learning}, 69(2--3), 169--192.}
{https://doi.org/10.1007/s10994-007-5016-8}
Το \emph{Online Newton Step} και το φράγμα $O(\log T)$ για exp-concave απώλειες.
Χρησιμοποιείται άμεσα στο τρίτο στάδιο OCO: η απώλεια ως προς τον πολλαπλασιαστή Kelly είναι
exp-concave, άρα το ONS είναι ο σωστός αλγόριθμος και όχι ο OGD. Αξίζει να προσεχθεί η μορφή
του σε μία διάσταση, όπου η προβολή στη νόρμα που ορίζει ο $A_t$ εκφυλίζεται σε απλή αποκοπή.

% =====================================================================
\section{Ειδικοί με δομή: specialists, tracking, εσωτερικό regret}
% =====================================================================

\src{freund1997}{Freund, Y., Schapire, R. E., Singer, Y., \& Warmuth, M. K. (1997). Using and
combining predictors that specialize. \emph{STOC}, 334--343.}
{https://doi.org/10.1145/258533.258616}
Η αναγωγή \emph{Sleeping Experts} / \emph{Specialists}, που είναι ο λόγος που η εργασία
μπορεί καν να δουλέψει με panel όπου η κάλυψη πηγαίνει από 7\% ως 100\%. Το κρίσιμο σημείο:
η εγγύηση regret κάθε expert μετριέται \emph{μόνο πάνω στους γύρους όπου ήταν ενεργός}, όχι
πάνω σε όλους. Χωρίς αυτό, ένας bookmaker που καλύπτει το 7\% των αγώνων θα αραίωνε την
εγγύηση όλων των υπολοίπων.

\src{herbster1998}{Herbster, M., \& Warmuth, M. K. (1998). Tracking the best expert.
\emph{Machine Learning}, 32(2), 151--178.}
{https://doi.org/10.1023/A:1007424614876}
Το Fixed-Share και η έννοια του \emph{regret παρακολούθησης}: σύγκριση όχι με τον καλύτερο
σταθερό expert αλλά με την καλύτερη \emph{ακολουθία} με φραγμένο πλήθος αλλαγών. Η
σημαντικότερη ιδέα για το δικό μας κεφάλαιο μαρκοβιανών αλυσίδων είναι ότι το Fixed-Share
\emph{είναι} η μπροστινή αναδρομή ενός κρυφού μαρκοβιανού μοντέλου με ομοιόμορφες
μεταβάσεις --- δηλαδή η εργασία περιείχε ήδη έναν αλγόριθμο με μαρκοβιανή αλυσίδα πριν
δοκιμαστεί ρητά οποιοσδήποτε άλλος.

\src{blum2007}{Blum, A., \& Mansour, Y. (2007). From external to internal regret.
\emph{Journal of Machine Learning Research}, 8, 1307--1324.}
{https://www.jmlr.org/papers/v8/blum07a.html}
Γενικεύει την αναγωγή των specialists και δίνει τη σύγχρονη, καθαρότερη διατύπωσή της.
Χρήσιμο κυρίως ως δεύτερη, πιο ευανάγνωστη πηγή για τους Sleeping Experts· η διάκριση
εξωτερικού και εσωτερικού regret δεν χρησιμοποιείται στην εργασία.

% =====================================================================
\section{Bandit ανάδραση}
% =====================================================================

\src{flaxman2005}{Flaxman, A. D., Kalai, A. T., \& McMahan, H. B. (2005). Online convex
optimization in the bandit setting: gradient descent without a gradient. \emph{SODA},
385--394.}
{https://dl.acm.org/doi/10.5555/1070432.1070486}
Η θεμελιώδης εργασία για bandit OCO και η πηγή του εκτιμητή \emph{ενός σημείου}. Το κεντρικό
κόλπο: μια διαταραχή $\delta u$ σε τυχαία κατεύθυνση επιτρέπει να κατασκευαστεί αμερόληπτη
εκτίμηση της κλίσης της \emph{εξομαλυμένης} συνάρτησης από μία μόνο βαθμωτή τιμή, με το
φράγμα να γίνεται $O(T^{3/4})$ αντί $O(T^{1/2})$. Αυτό που πρέπει να κρατηθεί για την
ερμηνεία των δικών μας αποτελεσμάτων είναι η \emph{διασπορά}: η εκτίμηση κλιμακώνεται ως
$d/\delta$, οπότε είναι αμερόληπτη αλλά σχεδόν μη ενημερωτική ανά γύρο.

\src{agarwal2010}{Agarwal, A., Dekel, O., \& Xiao, L. (2010). Optimal algorithms for online
convex optimization with multi-point bandit feedback. \emph{COLT}, 28--40.}
{https://www.learningtheory.org/colt2010/papers/032agarwal.pdf}
Ο εκτιμητής \emph{δύο σημείων} και η επαναφορά του φράγματος στο $O(\sqrt{T})$. Η λεπτομέρεια
που έχει σημασία για την τίμια ανάγνωση των αποτελεσμάτων μας: καθώς $\delta\to0$ η
πεπερασμένη διαφορά συγκλίνει στην ακριβή κατευθυντική παράγωγο, οπότε το μοντέλο δύο
ερωτημάτων παύει ουσιαστικά να είναι bandit για ομαλές συναρτήσεις. Το νούμερο του εκτιμητή
δύο σημείων διαβάζεται επομένως ως άνω φράγμα.

% =====================================================================
\section{Κανόνες βαθμολόγησης και βαθμονόμηση}
% =====================================================================

\src{gneiting2007}{Gneiting, T., \& Raftery, A. E. (2007). Strictly proper scoring rules,
prediction, and estimation. \emph{Journal of the American Statistical Association}, 102(477),
359--378.}
{https://doi.org/10.1198/016214506000001437}
Η πλήρης θεωρία των κατάλληλων κανόνων βαθμολόγησης. Το σημείο που χρησιμοποιείται: η
\emph{αναπαράσταση του Savage}, που δείχνει ότι κάθε κατάλληλος κανόνας αντιστοιχεί σε μια
κυρτή συνάρτηση πάνω στο simplex. Το log-loss και το Brier δεν είναι επομένως ανεξάρτητες
επιλογές αλλά δύο σημεία της ίδιας οικογένειας, κάτι που εξηγεί γιατί συμφωνούν σε όλες τις
ετυμηγορίες ενώ η ακρίβεια, που δεν ανήκει στην οικογένεια, δεν αναπαράγει καμία.

\src{murphy1973}{Murphy, A. H. (1973). A new vector partition of the probability score.
\emph{Journal of Applied Meteorology}, 12(4), 595--600.}
{https://doi.org/10.1175/1520-0450(1973)012\%3C0595:ANVPOT\%3E2.0.CO;2}
Η αποσύνθεση σε \emph{αξιοπιστία}, \emph{διακριτική ικανότητα} και \emph{αβεβαιότητα}. Αυτό
είναι που δίνει στο στάδιο βαθμονόμησης τη θεωρητική του δικαιολόγηση: ο όρος αξιοπιστίας
διορθώνεται χωρίς καμία επιπλέον πληροφορία, απλώς αναδιατάσσοντας μονότονα τις υπάρχουσες
πιθανότητες. Συνδέεται άμεσα με το εύρημα ότι η γραμμική συνάθροιση επιβαρύνει ακριβώς αυτόν
τον όρο.

\src{guo2017}{Guo, C., Pleiss, G., Sun, Y., \& Weinberger, K. Q. (2017). On calibration of
modern neural networks. \emph{ICML}.}
{https://arxiv.org/abs/1706.04599}
Η πηγή του \emph{temperature scaling} στη σύγχρονη μορφή του. Αυτό που χρειάζεται: ότι μία
μόνο παράμετρος αρκεί, ότι διατηρεί τη σειρά κατάταξης των ενδεχομένων --- άρα δεν αλλάζει
καμία απόφαση \en{argmax}, που εξηγεί γιατί βελτιώνει το log-loss αφήνοντας την ακρίβεια
αμετάβλητη --- και ο ορισμός του ECE που χρησιμοποιείται στη διαγνωστική ανάλυση.

% =====================================================================
\section{Αγορές στοιχηματισμού και μικροδομή}
% =====================================================================

\src{kyle1985}{Kyle, A. S. (1985). Continuous auctions and insider trading.
\emph{Econometrica}, 53(6), 1315--1335.}
{https://doi.org/10.2307/1913210}
Το υπόδειγμα όπου η τιμή κινείται γραμμικά με τη ροή εντολών και η ιδιωτική πληροφορία
ενσωματώνεται σταδιακά. Για εμάς είναι η θεωρητική βάση του γιατί η τιμή κλεισίματος είναι
καλύτερη πρόβλεψη από την τιμή ανοίγματος, και --- πιο συγκεκριμένα --- γιατί το πλεονέκτημά
της πρέπει να είναι μεγαλύτερο εκεί όπου η τιμή κινήθηκε περισσότερο. Αυτή είναι ποσοτική
πρόβλεψη που η εργασία ελέγχει και επιβεβαιώνει με διαβάθμιση 26 φορών.

\src{glosten1985}{Glosten, L. R., \& Milgrom, P. R. (1985). Bid, ask and transaction prices in
a specialist market with heterogeneously informed traders. \emph{Journal of Financial
Economics}, 14(1), 71--100.}
{https://doi.org/10.1016/0304-405X(85)90044-3}
Η \emph{δυσμενής επιλογή} ως εξήγηση του spread: ο διαμορφωτής χρεώνει περιθώριο επειδή το
γεγονός και μόνο ότι κάποιος δέχεται την τιμή του είναι κακή είδηση. Μεταφερόμενο στο δικό
μας πεδίο, αυτό εξηγεί τι \emph{είναι} το overround --- όχι τεχνικό εμπόδιο προς αφαίρεση
αλλά τίμημα προστασίας --- και γιατί το ύψος του διαφέρει ανά εταιρεία.

\src{levitt2004}{Levitt, S. D. (2004). Why are gambling markets organised so differently from
financial markets? \emph{The Economic Journal}, 114(495), 223--246.}
{https://doi.org/10.1111/j.1468-0297.2004.00226.x}
Τεκμηριώνει ότι ο bookmaker δεν εξισορροπεί απλώς τις δύο πλευρές αλλά \emph{λαμβάνει ενεργή
θέση}, εκμεταλλευόμενος συστηματικές προκαταλήψεις των παικτών. Αυτό είναι το επιχείρημα που
επιτρέπει να διαβαστούν οι διαφορές μεταξύ εταιρειών ως διαφορές \emph{επιχειρηματικού
μοντέλου} και όχι μόνο προγνωστικής ικανότητας --- και είναι η βάση της συζήτησης για το
γιατί ο ένας συγκεκριμένος οίκος δεν ξεπερνιέται.

\src{shin1992}{Shin, H. S. (1992). Prices of state contingent claims with insider traders, and
the favourite-longshot bias. \emph{The Economic Journal}, 102(411), 426--435.}
{https://doi.org/10.2307/2234478}
Το υπόδειγμα που παράγει το favorite--longshot bias \emph{ενδογενώς}, ως ορθολογική απόκριση
σε πληροφορημένους συναλλασσομένους, αντί να το αποδίδει σε ψυχολογική προκατάληψη. Η
σημασία για την εργασία είναι διπλή: δίνει εναλλακτική μέθοδο de-vig, και η παράμετρός του
$z$ έχει οικονομική ερμηνεία ως εκτιμώμενη αναλογία πληροφορημένων.

\src{shin1993}{Shin, H. S. (1993). Measuring the incidence of insider trading in a market for
state-contingent claims. \emph{The Economic Journal}, 103(420), 1141--1153.}
{https://doi.org/10.2307/2234552}
Η εφαρμόσιμη εκδοχή: πώς αντιστρέφεται το υπόδειγμα ώστε να ανακτηθούν δίκαιες πιθανότητες
και να εκτιμηθεί το $z$ ανά αγώνα. Από εδώ προέρχεται ο τύπος που υλοποιεί ο κώδικας. Το
$z$ χρησιμοποιείται επιπλέον ως δείκτης άφιξης πληροφορίας, όπου αποδεικνύεται πολύ
ασθενέστερος από την απευθείας μέτρηση της κίνησης της γραμμής.

\src{strumbelj2014}{Štrumbelj, E. (2014). On determining probability forecasts from betting
odds. \emph{International Journal of Forecasting}, 30(4), 934--943.}
{https://doi.org/10.1016/j.ijforecast.2014.01.001}
Εμπειρική σύγκριση των μεθόδων μετατροπής αποδόσεων σε πιθανότητες, με το συμπέρασμα ότι η
μέθοδος του Shin αποδίδει γενικά καλύτερα βαθμονομημένες πιθανότητες από την αναλογική
κανονικοποίηση. Είναι η δικαιολόγηση του να κατασκευαστεί ολόκληρο το σύνολο δεδομένων δύο
φορές και να ελεγχθεί κάθε κεντρικό αποτέλεσμα και με τις δύο.

% =====================================================================
\section{Μέγεθος στοιχήματος}
% =====================================================================

\src{kelly1956}{Kelly, J. L. (1956). A new interpretation of information rate. \emph{Bell
System Technical Journal}, 35(4), 917--926.}
{https://doi.org/10.1002/j.1538-7305.1956.tb03809.x}
Το κριτήριο Kelly, με την αρχική του διατύπωση μέσω ρυθμού πληροφορίας. Το ουσιώδες είναι ότι
η λογαριθμική χρησιμότητα δεν επιβάλλεται ως υπόθεση περί προτιμήσεων αλλά προκύπτει από την
\emph{επαναληπτική} δομή του προβλήματος: το κεφάλαιο πολλαπλασιάζεται, άρα το άθροισμα των
λογαρίθμων είναι το σωστό μέγεθος προς μεγιστοποίηση.

\src{breiman1961}{Breiman, L. (1961). Optimal gambling systems for favorable games.
\emph{Proceedings of the Fourth Berkeley Symposium on Mathematical Statistics and
Probability}, 1, 65--78.}
{https://projecteuclid.org/euclid.bsmsp/1200512363}
Η αυστηρή απόδειξη της ασυμπτωτικής βελτιστότητας: η στρατηγική Kelly μεγιστοποιεί τον ρυθμό
αύξησης σχεδόν βεβαίως και ελαχιστοποιεί τον αναμενόμενο χρόνο για την επίτευξη οποιουδήποτε
μεγάλου στόχου. Αυτό είναι το σημείο που χρειάζεται για να δικαιολογηθεί η επιλογή χωρίς
επίκληση σε υποκειμενική χρησιμότητα.

\src{cover1991}{Cover, T. M. (1991). Universal portfolios. \emph{Mathematical Finance}, 1(1),
1--29.}
{https://doi.org/10.1111/j.1467-9965.1991.tb00002.x}
Η online εκδοχή του προβλήματος κατανομής κεφαλαίου, με regret απέναντι στο καλύτερο σταθερό
αναπροσαρμοζόμενο χαρτοφυλάκιο. Χρήσιμο εδώ κυρίως \emph{αρνητικά}: εξηγεί γιατί η καθαρή
διατύπωση δεν μεταφέρεται, αφού ο comparator προϋποθέτει σταθερό σύνολο περιουσιακών
στοιχείων, ενώ εδώ τα «στοιχεία» είναι τα ενδεχόμενα ενός διαφορετικού αγώνα κάθε γύρο.

% =====================================================================
\section{Στατιστική και αριθμητικά εργαλεία}
% =====================================================================

\src{kunsch1989}{Künsch, H. R. (1989). The jackknife and the bootstrap for general stationary
observations. \emph{The Annals of Statistics}, 17(3), 1217--1241.}
{https://doi.org/10.1214/aos/1176347265}
Το \emph{moving block bootstrap}. Το σημείο που χρειάζεται: γιατί η δειγματοληψία μπλοκ
διατηρεί την τοπική εξάρτηση και γιατί ένα απλό $t$-test θα \emph{υποτιμούσε} τη διασπορά σε
θετικά συσχετισμένη ακολουθία, δηλαδή θα παρήγαγε ψευδώς σημαντικά αποτελέσματα. Χρησιμοποιείται
σε κάθε σύγκριση της εργασίας χωρίς εξαίρεση.

\src{hall1995}{Hall, P., Horowitz, J. L., \& Jing, B.-Y. (1995). On blocking rules for the
bootstrap with dependent data. \emph{Biometrika}, 82(3), 561--574.}
{https://doi.org/10.1093/biomet/82.3.561}
Η επιλογή του μήκους μπλοκ ως ανταλλαγή μεροληψίας--διασποράς, με αποτέλεσμα
$b \propto n^{1/3}$. Για το δικό μας $n$ δίνει $b\approx42$, δηλαδή περίπου μιάμιση
αγωνιστική. Το ίδιο μέγεθος επανεμφανίζεται αυτόνομα ως ορίζοντας εξάρτησης στη μελέτη του
βήματος, γεγονός που αξίζει να προσεχθεί γιατί προέκυψε από εντελώς διαφορετικό συλλογισμό.

\src{duchi2008}{Duchi, J., Shalev-Shwartz, S., Singer, Y., \& Chandra, T. (2008). Efficient
projections onto the $\ell_1$-ball for learning in high dimensions. \emph{ICML}.}
{https://doi.org/10.1145/1390156.1390191}
Ο αλγόριθμος προβολής στο simplex σε $O(n\log n)$. Χρειάζεται η κλειστή μορφή --- κατώφλι
$\theta$ και αποκοπή στο μηδέν --- γιατί εκτελείται δεκάδες χιλιάδες φορές και γιατί η
\emph{αραιοποιητική} της φύση έχει ουσιαστικές συνέπειες: μηδενίζει βάρη αντί να τα μικραίνει,
και η επαναλαμβανόμενη εφαρμογή της σε μεταβαλλόμενο σύνολο ενεργών experts παράγει κίνηση
βαρών ακόμη και χωρίς καμία μάθηση.

\src{duchi2011}{Duchi, J., Hazan, E., \& Singer, Y. (2011). Adaptive subgradient methods for
online learning and stochastic optimization. \emph{JMLR}, 12, 2121--2159.}
{https://www.jmlr.org/papers/v12/duchi11a.html}
Το AdaGrad. Στην εργασία εμφανίζεται ως δοκιμασμένη και απορριφθείσα εναλλακτική στον
χειροκίνητο συντονισμό του βήματος: ανακτά μικρό μόνο μέρος του οφέλους, και η εκδοχή ανά
expert αποδίδει χειρότερα από το να μην γίνει τίποτα.

% =====================================================================
\section{Μαρκοβιανές αλυσίδες}
% =====================================================================

\src{sun2018}{Sun, T., Sun, Y., \& Yin, W. (2018). On Markov chain gradient descent.
\emph{NeurIPS}.}
{https://arxiv.org/abs/1809.04216}
\emph{Δεν} αφορά μετάβαση βαρών, παρά το όνομα: μελετά SGD όπου τα δείγματα προέρχονται από
τροχιά μαρκοβιανής αλυσίδας αντί να είναι ανεξάρτητα. Η πρακτική συνέπεια που αξιοποιεί η
εργασία: αν ο ορίζοντας εξάρτησης είναι $\tau$, το ενεργό πλήθος δειγμάτων είναι $\sim t/\tau$
και το βήμα δικαιολογείται διογκωμένο κατά $\sqrt{\tau}$. Αυτό δίνει μερική θεωρητική εξήγηση
του εμπειρικά συντονισμένου $\times100$, χωρίς καμία σάρωση πλέγματος.

\vspace{1.5em}
\hrule
\vspace{0.8em}
{\small
\textbf{Σημείωση για τους συνδέσμους.} Οι σύνδεσμοι δόθηκαν από γνώση και όχι από αυτόματη
επαλήθευση ένας προς έναν. Τα DOI των περιοδικών άρθρων είναι σταθερά και υψηλής
βεβαιότητας· οι σύνδεσμοι προς πρακτικά συνεδρίων (ACM, COLT, AAAI) είναι οι πιθανότεροι να
χρειαστούν διόρθωση. Συνιστάται δειγματοληπτικός έλεγχος πριν την κατάθεση.
}

\end{document}
